\section{Grouped laws for exact calculation}\label{app:groups}
The following formulas specify a small sufficient-statistic experiment.
They are implemented with rational arithmetic in \texttt{code/cycle\_laws.py}.
Let $z$ be the number of zero labels, and set $p_0=p_N(0)$,
$p_1=p_N(a)$ for any $a\ne0$.

For odd $\ell=2r+1$, the independent-model mass of the group $z$ is
\[
Q_z=\binom{\ell}{z}p_0^z((q-1)p_1)^{\ell-z},
\]
and its likelihood ratio is
\[
L_z=1+N^{-r}
\left(\frac{(q-1)N}{N-q}\right)^z(-1)^{\ell-z}.
\]
Its symplectic mass is $Q_zL_z$.

For even $\ell=2k$, each group with $z\ge1$ has the same $Q_z$ and
\[
L_z
=1+\frac{N^k(N-q)(q-1)^{z-1}(-1)^z}
{(N-q)^z N^{\ell-z}}.
\]
When $z=0$, split into resonant and nonresonant groups:
\[
Q_{\rm res}=(q-1)^{\ell-1}p_1^\ell,\qquad
Q_{\rm non}=(q-2)(q-1)^{\ell-1}p_1^\ell,
\]
\[
L_{\rm res}
=1+N^{-k}\frac{N-q+(q-2)(N-1)q^k}{q-1},\qquad
L_{\rm non}
=1+N^{-k}\frac{N-q-(N-1)q^k}{q-1}.
\]
Omit the nonresonant group for $q=2$. Some symplectic masses can vanish for
small parameters; the formulas remain valid.

\section{Proof dependencies and formalization boundary}
The argument has the following dependencies:
\begin{enumerate}
\item Nondegeneracy of the alternating form gives the one-edge law.
\item Leaf removal gives forest independence.
\item Additive-character orthogonality gives $F_sF_t=NP_{s/t}-J$.
\item The cyclic trace gives the complete cycle Fourier law.
\item Additive inversion and one elementary Gauss-sum identity give the
pointwise cycle law.
\item Finite sums give exact total variation and chi-square values.
\item Symplectic transitivity and the block-cut tree give cactus factorization.
\item Taylor expansion, independence, and a central limit theorem give the
critical profile.
\item Conditioning on independent anchors gives the complete-bipartite
comparison, and the same product lemma gives its distinct critical scale.
\end{enumerate}
No step imports RH, GRH, an unproved finite-field estimate, or a numerical
extrapolation. A Lean formalization would need to implement the finite
probability spaces, the character matrices and their invariant subspaces,
the graph decomposition, and the triangular-array limit theorem.
The accompanying \texttt{FORMALIZATION.txt} gives more precise interfaces.

